% ===========================================================================
% Poste 13 — Aide de dernier recours
% ===========================================================================
\poste{13}{Aide de dernier recours (aide sociale)}{QC\_adr}

\pdesc Prestation mensuelle du programme d'aide sociale, versée aux
personnes et familles sans ressources suffisantes. C'est un
\emph{poste-intrant} : son montant entre dans le revenu familial net et
l'AFNI (équations~\ref{eq:rfn} et~\ref{eq:afni}) et déclenche l'exonération
de la prime RAMQ --- l'aide sociale alimente donc plusieurs autres postes.

\pobj Garantir un revenu minimal de subsistance (couverture des besoins de
base), tout en préservant une incitation au travail par l'exemption
partielle des gains.

\pregle Le calcul est mensuel, puis annualisé ($\times 12$). La prestation
est \emph{nulle dès qu'un adulte du ménage a 65~ans ou plus} (le système
bascule alors vers la Sécurité de la vieillesse, poste~17).

\begin{enumerate}
  \item \emph{Prestation de base mensuelle} : montant selon la composition
        ($b_{seul}$ ou $b_{couple}$), plus \emph{un seul} des ajustements
        suivants (par ordre de priorité) :
        \begin{itemize}
          \item les deux conjoints ont 58~ans et plus : $+\D{285}$ (2025) ;
          \item un adulte a 58~ans et plus, \emph{ou} le ménage est une
                famille monoparentale avec un enfant de moins de 5~ans
                (contrainte temporaire à l'emploi) : $+\D{166}$ (2025) ;
          \item adulte seul de moins de 50~ans sans enfant : $+\D{50}$.
        \end{itemize}
  \item \emph{Revenu compté}. Les \emph{gains de travail}, nets des
        cotisations (RRQ, RQAP, AE), bénéficient d'une exemption mensuelle
        $x$ (\D{200} seul, \D{300} couple) :
        $g = \pos{R_{net}/12 - x}$. Les revenus \emph{autres} --- la
        pension d'un retraité de moins de 65~ans, par exemple --- sont
        comptés \emph{au complet}, sans exemption : $c = g + A/12$ ;
  \item \emph{Prestation nette} : $p = \pos{b - c}$ ;
  \item \emph{Incitation au travail} : si la prestation nette est positive,
        \pc{25} des gains de travail comptés sont remboursés :
        $i = 0{,}25 \times g$ si $p > 0$, sinon $0$ ;
  \item \emph{Montant annuel} : $QC_{adr} = 12 \times (p + i)$.
\end{enumerate}

L'effet de l'incitation est visible dans le graphique des taux marginaux de
l'application : sur la plage où l'aide se retire, chaque dollar de travail
net ne réduit la prestation que de 75~cents.

\pparams

\begin{center}
\small
\begin{tabular}{l r r}
\toprule
Paramètre (mensuel) & 2025 & 2026 \\
\midrule
Prestation de base --- 1 adulte ($b_{seul}$)    & \D{829}  & \D{845} \\
Prestation de base --- couple ($b_{couple}$)    & \D{1258} & \D{1283} \\
Ajustement --- un adulte 58 ans et + (ou contrainte temporaire) & \D{166} & \D{169} \\
Ajustement --- deux adultes 58 ans et +         & \D{285}  & \D{291} \\
Ajustement --- adulte seul de moins de 50 ans   & \D{50}   & \D{50} \\
Exemption de gains --- 1 adulte ($x$)           & \D{200}  & \D{200} \\
Exemption de gains --- couple ($x$)             & \D{300}  & \D{300} \\
Incitation au travail (part des gains exemptée) & \pc{25}  & \pc{25} \\
\bottomrule
\end{tabular}
\end{center}

% ============================================================================
% GÉNÉRÉ par scripts/exemples-guide.ts — NE PAS ÉDITER À LA MAIN.
% Régénérer : npm run exemples (chaque chiffre est vérifié contre le moteur).
% ============================================================================
\pexemples Calcul \emph{mensuel} puis annualisé : prestation de base (selon la
composition et l'âge), moins le revenu de travail \emph{net des cotisations}
au-delà de l'exemption mensuelle (\D{200} pour un adulte seul,
\D{300} pour un couple) ; si la prestation reste positive,
\pc{25} des gains comptés sont remboursés (incitation au
travail).

\medskip\noindent\textbf{M1 --- personne vivant seule, 25~ans, aucun revenu.}
Aucun revenu. Adulte seul de moins de 50~ans sans enfant : ajustement
de \D{50} en sus de la base de \D{829}.
\begin{align*}
\text{prestation mensuelle} &= \num{829} + \num{50} = \num{879} \\
\text{aide annuelle} &= \num{879} \times 12 = \num{10548}
\end{align*}
Aide de dernier recours : \D{10548}.

\medskip\noindent\textbf{M2 --- personne vivant seule, 30~ans, revenu de travail de \D{9000}.}
Revenu de travail de \D{9000}, net des cotisations (postes 1 à 3)
\D{514.36}, soit \D{8485.64} par année.
\begin{align*}
\text{gains comptés} &= \pos{\num{8485.64} \div 12 - \num{200}} = \num{507.14} \\
\text{prestation nette} &= \pos{\num{879} - \num{507.14}} = \num{371.86} \\
\text{incitation} &= \pc{25} \times \num{507.14} = \num{126.78} \\
\text{aide annuelle} &= (\num{371.86} + \num{126.78}) \times 12 = \num{5983.77}
\end{align*}
Aide de dernier recours : \D{5983.77}. Noter le rôle de l'exemption de
\D{200} par mois et de l'incitation au travail de
\pc{25}.

\medskip\noindent\textbf{M6 --- famille monoparentale, 35~ans, revenu de travail de \D{35000}, un enfant de 3~ans en garderie subventionnée (\D{2000} payés dans l'année).}
Parent seul d'un enfant de moins de 5~ans : la prestation de base
serait majorée de l'ajustement « contrainte temporaire à l'emploi »
(\D{166}). Mais les gains comptés
($\approx \num{2496.05}$
par mois) dépassent largement cette base : prestation nulle, donc pas
d'incitation non plus — aide de \D{0}.

\medskip\noindent\textbf{M7 --- couple, 45 et 44~ans, revenus de travail de \D{30000} et \D{0}, sans enfant.}
Couple (base de \D{1258}, exemption de \D{300}) :
les gains de travail comptés excèdent la base — aide nulle. L'aide de dernier
recours s'éteint vite : c'est un programme de \emph{dernier} recours.


\prefs Loi sur l'aide aux personnes et aux familles (RLRQ, c.~A-13.1.1) et
son règlement ; ministère de l'Emploi et de la Solidarité sociale
(barèmes publiés sur quebec.ca).
